%!TEX root = ../main.tex
% APÉNDICE A — Encuesta de validación

\chapter{Encuesta de validación}\label{apx:encuesta}

\begin{guia}
Los apéndices contienen material producido por el autor que respalda
el cuerpo pero lo interrumpiría. Acá va la encuesta completa: para qué
se hizo, cómo se diseñó y distribuyó, a cuántas personas llegó y el
cuestionario tal como lo vieron quienes respondieron. Los resultados
se analizan en el capítulo~\ref{cap:resultados}.
\end{guia}

\section{Introducción}

El objetivo de la encuesta es relevar cómo organizan sus eventos
potenciales usuarios de la plataforma (personas que ya hayan organizado
un casamiento, cumpleaños, fiesta de 15, evento empresarial u otro
evento similar), identificar los puntos más difíciles del proceso actual
de búsqueda y contratación de proveedores, y validar el interés en una
plataforma integral con marketplace, recomendación por \gls{ia} y
simulación del evento. Los resultados alimentan directamente el
capítulo~\ref{cap:marco-teorico} (estado de la problemática) y el
capítulo~\ref{cap:economica} (estimación de demanda).

\section{Metodología}

El cuestionario se armó en Google Forms, con 10 preguntas (opción única,
opción múltiple y escala 1--5, sin preguntas abiertas), y se distribuyó
de forma no probabilística por redes sociales y contactos personales del
autor. Estuvo abierto entre el 22/09/2026 y el 29/09/2026 y recibió 45
respuestas.

\begin{nota}
El muestreo es por conveniencia, no aleatorio, por lo que los resultados
no son representativos de toda la población y sirven como una primera
validación exploratoria, no como una medición estadísticamente
generalizable. Esta limitación se retoma en el
capítulo~\ref{cap:conclusiones}.
\end{nota}

\section{Presentación de la encuesta}

\completar{Texto introductorio exacto que vieron las personas
participantes (propósito, anonimato, duración estimada), tal como
apareció en la portada del formulario de Google Forms.}

\section{Cuestionario}

\begin{enumerate}
    \item ¿Organizaste algún evento (casamiento, cumpleaños, fiesta de
          15, evento empresarial, etc.)? (opción única: Sí / No)
    \item ¿Cuántas personas asistieron aproximadamente al evento? (opción
          única: Menos de 50 / 50 a 100 / 101 a 200 / Más de 200)
    \item ¿Qué medios usaste para buscar y contratar proveedores? (opción
          múltiple: WhatsApp / Instagram o Facebook / Google /
          Recomendación de conocidos / Alguna app o sitio tipo
          marketplace (BIG Fiestas, Casamientos.com.ar, etc.) / Contraté
          un organizador o wedding planner / Otro)
    \item Aproximadamente, ¿cuánto tiempo te llevó buscar y comparar
          proveedores? (opción única: Menos de 1 semana / Entre 1 y 4
          semanas / Más de 1 mes)
    \item ¿Cuántos proveedores distintos tuviste que contactar en total
          (salón, catering, fotografía, decoración, etc.)? (opción
          única: 1 a 3 / 4 a 6 / 7 o más)
    \item ¿Qué fue lo más difícil de organizar? (opción múltiple:
          Encontrar opciones disponibles / Comparar precios / Coordinar
          horarios entre proveedores / Llevar el control de pagos y
          presupuesto / Otro)
    \item Del 1 al 5, ¿qué tan complicado te resultó organizar el evento
          en general? (escala 1--5)
    \item Si existiera una sola plataforma donde buscar, comparar y
          contratar todos los servicios para tu evento, ¿la usarías?
          (opción única: Sí, seguro / Probablemente / No lo sé / No)
    \item ¿Te resultaría útil que la plataforma, usando \gls{ia}, te
          recomiende automáticamente proveedores según tu presupuesto,
          ubicación y cantidad de invitados? (opción única: Sí / No sé /
          No)
    \item ¿Te resultaría útil que la plataforma te permita simular el
          evento, por ejemplo viendo cómo cambian el costo y la
          organización si aumenta la cantidad de invitados? (opción
          única: Sí / No)
\end{enumerate}
